\documentclass[12pt,a4paper,oneside]{article}
\usepackage[brazil]{babel}
\usepackage[utf8]{inputenc}
\usepackage{olymp}
\usepackage{booktabs}
\usepackage{wrapfig}

\setcounter{problem}{0}

\begin{document}

\begin{problem}{Peso em outros planetas}{peso}{2}

\begin{wrapfigure}{r}{0.3\textwidth}
    \includegraphics[width=0.3\textwidth]{figuras/Earthset.png}
    \centering {\footnotesize Pôr-da-terra: foto do comandante\\da missão Artemis II}
\end{wrapfigure}

Na sexta-feira passada, os 4 astronautas da missão Artemis 2, da Nasa, retornaram com sucesso à Terra.
O ponto alto da missão foi um sobrevoo ao redor da Lua, permitindo que os astronautas vissem com seus próprios olhos o lado oculto da Lua.
Eles ainda viram um eclipse solar e um ``Pôr-da-Terra'' (foto ao lado).

Segundo a Nasa, esse foi um voo de teste para uma missão futura que deve levar humanos de volta à Lua e também para validar tecnologias essenciais para viagens mais longas,
por exemplo verificar o quão rapidamente os astronautas conseguem se ajustar a uma mudança de gravidade para o pouso na Lua e uma possível missão a Marte.
Devido à diferença de gravidade, o peso de uma pessoa na Lua é cerca de $1/6$ do seu peso na Terra, enquanto em Marte cai para quase $1/3$.

A tabela a seguir mostra valores mais precisos, que devem ser usados nesse exercício.

\begin{center}
\begin{tabular}{lr}
        \toprule
        \textbf{Corpo Celeste} & \textbf{Valor (\%)} \\
        \midrule
        Lua & 16.5\% \\
        Marte & 38\% \\
        Venus & 91\% \\
        Jupiter & 234\% \\
        \bottomrule
\end{tabular}
\end{center}

%Assim, uma pessoa de $70$ Kg na Terra pesa apenas $70 \times 0.165 = 11.55$ Kg na Lua, $70 \times 0.38 = 26.6$ Kg em Marte e impressionantes mais de $150$ em Júpiter.

%\Notes
%
%Observações, se tiver.

\InputFile

A entrada contém apenas uma linha, com um valor inteiro $K$, o peso em Kg de uma pessoa na Terra, e uma letra maiúscula $P$ informando um local.
Restrições: $50 \leq K \leq 120$ e $P$ pode ser `\texttt{L}', `\texttt{M}', `\texttt{V}' ou `\texttt{J}', respectivamente os corpos celestes na tabela.

\OutputFile

Escreva o peso da pessoa no local solicitado, formatado para uma casa decimal.

%\Notes

\Example

\begin{example}
\exmp{
70 L
}{
11.6
}%
\end{example}

\begin{example}
\exmp{
70 M
}{
26.6
}%
\end{example}

\begin{example}
\exmp{
100 L
}{
16.5
}%
\end{example}

\begin{example}
\exmp{
50 J
}{
117.0
}%
\end{example}

%Comentários sobre o exemplo, se necessário.
\small
Primeiro exemplo: uma pessoa de $70$ Kg na Terra pesará $70 \times 0.165 = 11.55 \approx 11.6$ Kg na Lua.\\
Último exemplo: uma pessoa de $50$ Kg na Terra pesará impressionantes $117.0$ Kg em Júpiter.

\end{problem}

\end{document}
